
\chapter{Local comparisons on the groups}

\section{Introduction}
We prove the identities necessary to compare the
invariant trace formula to its metaplectic analog in chapter 10.

\section{Weighted orbital integrals}

Certain weighted orbital integrals occur naturally in the invariant trace
formula. We will need some of their properties to put the trace formula in
stable form as well as to compare trace formulas.  

If $M\in \{A,G\}$, $f\in C_c^\infty (\overline{G(F)})_{gen}$, and $g\in 
\overline{M(F)}$ is regular, let
\[
J_M(g,f)=|D(g)|^{1/2}
\int_{\overline{G(F)}_g\backslash \overline{G(F)}}
f(x^{-1}gx)v_M(x)dx.
\]
If $F={\Bbb{Q}}_p$ or $F={\Bbb{R}}$ then
\[
v_M(x)=\left\{ 
\begin{array}{c}
1,\;\;\;\;M=G, \\ 
\log \sqrt{|1+n^2|_F},\;\;M=A,
\end{array}
\right. 
\]
where $x=a\left( 
\begin{array}{cc}
1 & n \\ 
0 & 1
\end{array}
\right) k$. We will need to compute the Fourier transform of $J_M(a,f)$ as a
function on $A$.

First, consider $M=G$.

\begin{lemma}
Let $f\in C_c^\infty (\overline{G(F)})_{gen}$. The function $m\longmapsto
J_G(m,f)$ from $\overline{A(F)}\rightarrow {\Bbb C}$ vanishes unless $m\in 
\overline{A(F)^N}$.
\end{lemma}

proof: The usual trick: Let $g=\left( 
\begin{array}{cc}
a & 0 \\ 
0 & a^{-1}
\end{array}
\right) $, $h=\left( 
\begin{array}{cc}
b & 0 \\ 
0 & b^{-1}
\end{array}
\right) $, so
\[
(g,1)^{-1}(h,1)(g,1)=(h,(a,b)^{-2}).
\]
Since the measure $dx$ is translation invariant, if $m=(h,1)$ then we have
\[
\int_{\overline{G(F)}_m\backslash \overline{G(F)}}f(x^{-1}mx)dx=(a,b)^{-2}%
\int_{\overline{G(F)}_m\backslash \overline{G(F)}}f(x^{-1}mx)dx.
\]
The function $b\longmapsto (a,b)^{-2}$ is trivial if and only if $a\in A(F)^N
$. $\Box $

As a corollary we shall want later,

\begin{lemma}
Let $f\in C_c^\infty (\overline{G(F)})_{gen}$. The function 
\[
m\longmapsto \int_{N(F)}\int_{\overline{K}}f(k^{-1}n^{-1}mnk)dkdn
\]
from $\overline{A(F)}\rightarrow {\Bbb C}$ vanishes unless
$m\in \overline{A(F)^N}$.
\end{lemma}

Since the weight $v_M(x)$ satisfies $v_M(ax)=v_M(x)$ for
all $a\in\overline{A(F)}$, the same proof yields the analogous
fact for the weighted orbital integrals:

\begin{lemma}
Let $f\in C_c^\infty (\overline{G(F)})_{gen}$. The function
$m\longmapsto J_M(m,f)$ from $\overline{A(F)}\rightarrow {\Bbb C}$
vanishes unless $m\in \overline{A(F)^N}$.
\end{lemma}

Indeed,
\begin{lemma}
Let $f\in C_c^\infty (\overline{G(F)})_{gen}$. The function 
\[
m\longmapsto
\int_{N(F)}\int_{\overline{K}}f(k^{-1}n^{-1}mnk)v_M(nk)dkdn
\]
from $\overline{A(F)}\rightarrow {\Bbb C}$ vanishes unless
$m\in \overline{A(F)^N}$.
\end{lemma}

Note that the centralizer in $\overline{G(F)}$ of a regular element
in $a\in A(F)^N$ is $\overline{A(F)}$. As a consequence of the
above lemmas, we may therefore write,
for a regular $a\in \overline{A(F)}$,
\begin{equation}
J_M(a,f)=
\left\{
\begin{array}{c}
|D(g)|^{1/2}\int_{\overline{A(F)}\backslash \overline{G(F)}}
f(x^{-1}gx)v_M(x)dx,\ \ \ a\in \overline{A(F)^N},\\
0,\ \ a\in \overline{A(F)}- \overline{A(F)^N}.
\end{array}
\right.
\end{equation}

\subsection{Elliptic orbital integrals}

We now discuss the elliptic orbital integrals $J_G(\gamma)$,
$\gamma\in \overline{T}_{reg}$, where $\overline{T}_{reg}$
denotes the subset of elements of the form $(x,\zeta)$
where $x$ has distinct eigenvalues.

By Lemma 3(h), we have
\[
\overline{G(F)}_\gamma
=\overline{T},\ \ \ \ \ \
\gamma\in \overline{T}_{reg}.
\]
We have
\begin{equation}
\begin{array}{c}
\int_{\overline{G(F)}_\gamma\backslash \overline{G(F)}}
f(x^{-1}\gamma x){dx\over dt}
=\int_{\overline{T}\backslash \overline{G(F)}}
f(x^{-1}\gamma x){dx\over dt}\\
=\int_{T\backslash G(F)}
f(x^{-1}gx,1){dx\over dt},
\end{array}
\end{equation}
where $\gamma =(g,1)\in \overline{T}_{reg}$.

\subsection{Hyperbolic orbital integrals (n=2)}

If $n=2$ then by Lemma 3(g), we have
\[
\overline{G(F)}_\gamma
=\overline{A(F)},\ \ \ \ \ \
\gamma\in \overline{A(F)}_{reg}.
\]
We have
\begin{equation}
\begin{array}{c}
\int_{\overline{G(F)}_\gamma\backslash \overline{G(F)}}
f(x^{-1}\gamma x){dx\over da}
=\int_{\overline{A(F)}\backslash \overline{G(F)}}
f(x^{-1}\gamma x){dx\over da}\\
=\int_{A(F)\backslash G(F)}
f(x^{-1}gx,1){dx\over da},
\end{array}
\end{equation}
where $\gamma =(g,1)\in \overline{A(F)}_{reg}$.

This equation suggests, somewhat misleadingly, a ``naive'' matching between
smooth functions: simply send $\overline{f}(g,1)$ to
$f=\overline{f}(g,1)$. The correct matching will be defined in the
next section.

\subsection{Unipotent orbital integrals}

Using (2.6) a straightforward calculation involving only the definitions
yields the following result:
\[
\overline{G(F)}_\gamma
=\overline{N(F)},\ \ \ \ \ \
\gamma=
(
\pm \left(
\begin{array}{cc}
1 & 1\\
0 & 1
\end{array}
\right)
,1)\in \overline{N(F)}.
\]
We have
\begin{equation}
\begin{array}{c}
\int_{\overline{G(F)}_\gamma\backslash \overline{G(F)}}
f(x^{-1}\gamma x){dx\over dn}
=\int_{\overline{N(F)}\backslash \overline{G(F)}}
f(x^{-1}\gamma x){dx\over dn}\\
=\int_{N(F)\backslash G(F)}
f(x^{-1}gx,1){dx\over dn},
\end{array}
\end{equation}
where $\gamma $ is as above.

\subsection{Weighted orbital integrals (n=2)}

Assume $n=2$, so the two-fold cover of $A(F)$ is abelian.
If $\gamma =(g,1)\in \overline{A(F)}_{reg}$ and if
\[
g=
\pm \left(
\begin{array}{cc}
a & 0\\
0 & a^{-1}
\end{array}
\right)
\in A(F),
\]
then 
\begin{equation}
\begin{array}{c}
\int_{\overline{G(F)}_\gamma\backslash \overline{G(F)}}
f(x^{-1}\gamma x)v_A(x){dx\over da}
=\int_{\overline{A(F)}\backslash \overline{G(F)}}
f(x^{-1}\gamma x)v_A(x){dx\over da}\\
=\int_{A(F)\backslash G(F)}
f(x^{-1}gx,1)v_A(x){dx\over da},
\end{array}
\end{equation}
as in \S 5.2.2.

If
\[
\gamma=
(
\pm \left(
\begin{array}{cc}
1 & 0\\
0 & 1
\end{array}
\right)
,1)\in Z_{\overline{G}}(F),
\]
we have
\begin{equation}
\begin{array}{c}
J_M(\gamma,f)=
\lim\Sb a\rightarrow 1\\ a\ {\rm regular}\endSb
\int_{\overline{G(F)}_{a\gamma}\backslash \overline{G(F)}}
f(x^{-1}\gamma a x)v_A(x){dx\over dt}\\
=\lim\Sb a\rightarrow 1\\ a\ {\rm regular}\endSb
\int_{\overline{A(F)}\backslash \overline{G(F)}}
f(x^{-1}\gamma a x)v_A(x){dx\over dt}\\
=\lim\Sb a\rightarrow 1\\ a\ {\rm regular}\endSb
\int_{A(F)\backslash G(F)}
f(x^{-1}p(\gamma)a x,1)v_A(x){dx\over dt},
\end{array}
\end{equation}

\section{Corresponding functions, $n=2$}

We define the notion of a ``smooth matching''.

\subsection{Corresponding functions on $A$}

We concentrate first on
the set $A(F)_{reg}$.

\subsubsection{Characters of $\overline{A(F)}$ for $n=2$}

We assume $p>2$ throughout this subsection unless otherwise
stated,

Let
\[
A_2(F)=\{
a=\left(
\begin{array}{cc}
x & 0\\
0 & x^{-1}
\end{array}
\right)
\ |\ \nu (x)\ {\rm \ even}\},
\]
so $\overline{A_2(F)}
=A_2(F)\times \mu_2$ (direct product).
The elements of $\Pi(\overline{A(F)})$ are obtained by extending
the elements of $\Pi(\overline{A_2(F)})$.
If $\chi_1\in \Pi(\overline{A_2(F)})$ then
$\chi_1=\chi_2\times id^{-1}$,
for some $\chi_2\in \Pi(A_2(F))$

We extend $\chi_1$ to $\overline{A(F)}$ by
\begin{equation}
\chi(a,\zeta)=\chi_\epsilon (a,\zeta)=
\left\{
\begin{array}{c}
\zeta^{-1}\chi_1(x),\ \ \ a\in A_2(F),\\
\epsilon\zeta^{-1}(x,\pi)^{-1}\chi_1(x/\pi),
\ \ \ a\in A(F)-A_2(F),
\end{array}
\right.
\end{equation}
where $\epsilon^2 = (\pi,\pi)^{-1}\chi_1(\pi^2)$. We claim
that this defines a character. (This is left to
the reader but see also \cite{GS}, Lemma 2.1 for an equivalent form.)
Therefore, the characters of $\overline{A(F)}$ are in 2-1 correspondence
with the characters of $A_2(F)$.

Since $A(F)/A_2(F)$ is order 2, the characters of $A(F)$ are
also in 2-1 correspondence with the characters of $A_2(F)$.
Explicitly, given any $\chi_1\in \Pi(A_2(F))$, we
may extend $\chi_1$ to $A(F)$ by
\begin{equation}
\chi(a)=\chi_\theta (a)=
\left\{
\begin{array}{c}
\chi_1(x),\ \ \ a\in A_2(F),\\
\theta \chi_1(x/\pi), \ \ \ a\in A(F)-A_2(F),
\end{array}
\right.
\end{equation}
where $\theta^2 = \chi_1(\pi^2)$.

\begin{definition} Let $\chi_1\in \Pi(A_2(F))$ be given.
We say that
$\chi_\epsilon\in \Pi(\overline{A(F)})$ {\bf corresponds}
to $\chi_\theta \in \Pi(A(F))$ if $\epsilon =\theta$.
If $\chi \in \Pi(A(F))$ is given we let
$\overline{\chi}\in \Pi(\overline{A(F)})$ denote
its corresponding genuine character.
\end{definition}

\subsubsection{Corresponding functions on $\overline{A(F)}_{reg}$}

First, we investigate the conditions on
$\phi \in C_c^\infty(A(F))$,
$\overline{\phi}
\in C_c^\infty(\overline{A(F)})_{gen}$
under which
\[
\phi^{\wedge}(\chi)=
\overline{\phi}^{\wedge}(\overline{\chi}),
\]
if $\chi \in \Pi(A(F))$.

We have
\begin{equation}
\begin{array}{c}
\phi^{\wedge}(\chi)=
\int_{A(F)}\phi(a)\chi(a) da\\
=\int_{A_2(F)}\phi(a)\chi_1(x_a) da +
\epsilon \int_{A(F)-A_2(F)}\phi(a)\chi_1(x_a/\pi) da
\end{array}
\end{equation}
where
\[
a=\left(
\begin{array}{cc}
x_a & 0\\
0 & x_a^{-1}
\end{array}
\right).
\]

Similarly, 
\begin{equation}
\begin{array}{c}
\overline{\phi}^{\wedge}(\overline{\chi})=
\int_{A(F)}\overline{\phi}(a,1)\overline{\chi}(a,1) da\\
=\int_{A_2(F)}\overline{\phi}(a,1)\chi_1(x_a) da +
\epsilon \int_{A(F)-A_2(F)}\overline{\phi}(a,1)\chi_1(x_a/\pi) da.
\end{array}
\end{equation}

\begin{definition} We say that
$\phi \in C_c^\infty(A(F))$ {\bf corresponds} to
$\overline{\phi}
\in C_c^\infty(\overline{A(F)})_{gen}$
if
\[
\overline{\phi}(a,1)=
\left\{
\begin{array}{c}
f(a),\ \ \ \ a\in A_2(F),\\
(x_a,\pi)^{-1}f(a),\ \ \ \ a\in A(F)-A_2(F).
\end{array}
\right.
\]
\end{definition}

We have proven

\begin{lemma} If
$\phi \in C_c^\infty(A(F))$ corresponds to
$\overline{\phi}\in C_c^\infty(\overline{A(F)})_{gen}$
then
\[
\phi^{\wedge}(\chi)=
\overline{\phi}^{\wedge}(\overline{\chi}),
\]
if $\chi \in \Pi(A(F))$.
\end{lemma}

\begin{definition} We say that
$f \in C_c^\infty(G(F))$ {\bf corresponds on $A(F)_{reg}$} to
$\overline{f}\in C_c^\infty(\overline{G(F)})_{gen}$
if $F_f\in C_c^\infty(A(F))$ corresponds to
$F_{\overline{f}}\in C_c^\infty(\overline{A(F)})_{gen}$
in the sense above.
\end{definition}

The following remark holds even if $p=2$:
For $f$ of sufficiently small support, it can be
shown by a direct calculation of the orbital integrals that if
$f \in C_c^\infty(G(F))$ corresponds on $A(F)_{reg}$ to
$\overline{f}\in C_c^\infty(\overline{G(F)})_{gen}$
then the weighted orbital integral of
$f$ corresponds to the weighted orbital integral of
$\overline{f}$ in the sense above.
This remark shall be used later in chapter 10.

\subsubsection{Case $p=2$}

Now we discuss the case $p=2$ briefly. Let
\[
A_2'(F)=\{
a=\left(
\begin{array}{cc}
x & 0\\
0 & x^{-1}
\end{array}
\right)
\ |\ \nu (x)\ {\rm \ even\ and\ }x/\pi^{-\nu(x)}\in B\},
\]
where $B\subset {\cal{O}}_F$ is choosen to be finite index and
so that the Hilbert symbol on $B\times B$ is trivial.
Thus $\overline{A_2'(F)}
=A_2'(F)\times \mu_2$ (direct product).
The elements of $\Pi(\overline{A(F)})$ are obtained by extending
the elements of $\Pi(\overline{A_2'(F)})$.
If $\chi_1\in \Pi(\overline{A_2'(F)})$ then
$\chi_1=\chi_2\times id^{-1}$,
for some $\chi_2\in \Pi(A_2'(F))$.

The elements of $\Pi(A(F))$ are obtained by extending
the elements of $\Pi(A_2'(F))$. Since everything is abelian,
\begin{lemma} There exists a (non-canonical) 1-1 correspondence
between elements of $\Pi(\overline{A(F)})$ and
elements of $\Pi(A(F))$. 
\end{lemma}

Since there are several, fix one
once and for all.
If $\chi\in \Pi(A(F))$ then
we denote by $\overline{\chi}\in \Pi(\overline{A(F)})$
the corresponding character.

\begin{definition}
We say that $\phi \in C_c^\infty(A(F))$ {\bf corresponds} to
$\overline{\phi}\in C_c^\infty(\overline{A(F)})_{gen}$
if
\[
\phi^{\wedge}(\chi)=
\overline{\phi}^{\wedge}(\overline{\chi}),
\]
for all $\chi \in \Pi(A(F))$.
\end{definition}

In the case $p=2$ we repeat the above definition.

\begin{definition} We say that
$f \in C_c^\infty(G(F))$ {\bf corresponds on $A(F)_{reg}$} to
$\overline{f}\in C_c^\infty(\overline{G(F)})_{gen}$
if $F_f\in C_c^\infty(A(F))$ corresponds to
$F_{\overline{f}}\in C_c^\infty(\overline{A(F)})_{gen}$
in the sense above.
\end{definition}

\subsection{Corresponding functions on
$\overline{A({\Bbb{R}})}$, $n=2$}

Let $F={\Bbb{R}}$ in this subsection.
The genuine characters of $\overline{A(F)}$, $n=2$, are
described in \cite{GS} and they are shown to be in 1-1 correspondence with
the characters of $A(F)$.
If $\chi\in \Pi(A(F))$ then
we denote by $\overline{\chi}\in \Pi(\overline{A(F)})$
the corresponding character.

\begin{definition}
We say that $\phi \in C_c^\infty(A(F))$ {\bf corresponds} to
$\overline{\phi}\in C_c^\infty(\overline{A(F)})_{gen}$
if
\[
\phi^{\wedge}(\chi)=
\overline{\phi}^{\wedge}(\overline{\chi}),
\]
for all $\chi \in \Pi(A(F))$.
\end{definition}

\begin{definition} We say that
$f \in C_c^\infty(G(F))$ {\bf corresponds on $A(F)_{reg}$} to
$\overline{f}\in C_c^\infty(\overline{G(F)})_{gen}$
if $F_f\in C_c^\infty(A(F))$ corresponds to
$F_{\overline{f}}\in C_c^\infty(\overline{A(F)})_{gen}$
in the sense above.
\end{definition}

Note that as far as the comparison of orbital integrals is concerned,
the assumption that the functions be infinitely differentiable
may be weakened, if desired, to the assumption that the
functions be compactly supported and piecewise
continuous. The space of infinitely differentiable
functions of compact support is dense in the space of
piecewise continuous functions of compact support, so if we can
establish the comparison of orbital integrals for functions
in one space we may deduce it for functions in the
other space. Using an approximation argument involving
step functions, it can be
shown by a direct calculation of the orbital integrals that if
$f \in C_c^\infty(G({\Bbb{R}}))$ corresponds on
$A({\Bbb{R}})_{reg}$ to
$\overline{f}\in C_c^\infty(\overline{G({\Bbb{R}})})_{gen}$
then the weighted orbital integral of
$f$ corresponds to the weighted orbital integral of
$\overline{f}$ in the sense above.
This remark shall be used later in chapter 10.

We shall consider the correspondence between characters of
$A({\Bbb{R}})$ and its two-fold cover in a little more
detail.

Let $a(x)$ denote the diagonal matrix with entries $x,x^{-1}$.
Let
\[
(x,y)_2=
\left\{
\begin{array}{c}
1,\ \ \ \ {\rm if \ }x>0,\\
-1,\ \ \ \ {\rm if \ }x<0.
\end{array}
\right.
\]
Let $\chi_1$ be a character on ${\Bbb{R}}_+^\times$. Define
$\chi$ on $A({\Bbb{R}})$ by
\[
\chi(a(x))=\chi_\epsilon(a(x))
=
\left\{
\begin{array}{c}
\chi_1(x),\ \ \ \ {\rm if \ }x>0,\\
\epsilon \chi_1(-x),\ \ \ \ {\rm if \ }x<0,
\end{array}
\right.
\]
where $\epsilon =\pm 1$. Define
$\overline{\chi}$ on $\overline{A({\Bbb{R}})}$ by
\[
\overline{\chi}(a(x),\zeta)=\overline{\chi}_\eta(a(x),\zeta)
=
\left\{
\begin{array}{c}
\zeta^{-1}\chi_1(x),\ \ \ \ {\rm if \ }x>0,\\
\zeta^{-1}\eta \chi_1(-x),\ \ \ \ {\rm if \ }x<0,
\end{array}
\right.
\]
where $\eta=\pm i$. We say that
$\chi$ {\bf corresponds} to $\overline{\chi}$ if
$\eta =\epsilon i$.

The above definition boils down to saying that
a smooth function $f$ on $A({\Bbb{R}})$
{\bf corresponds to} a genuine smooth function
$\overline{f}$ on $\overline{A({\Bbb{R}})}$ if
\begin{equation}
f(a,1)=
\left\{
\begin{array}{c}
f(a),\ \ \ \ {\rm if \ }a=a(x),\ x>0,\\
\epsilon f(a),\ \ \ \ {\rm if \ }a=a(x),\ x<0,
\end{array}
\right.
\end{equation}

\subsection{Corresponding functions on the elliptic set}

In this subsection, $F$ may be either real or p-adic.

Finally, we concentrate on the set $T(F)_{reg}$, $T$ elliptic.

\begin{definition} We say that
$f \in C_c^\infty(G(F))$ {\bf corresponds on $T(F)_{reg}$} to
$\overline{f}\in C_c^\infty(\overline{G(F)})_{gen}$
if the orbital integral $F_f\in C_c^\infty(T(F)_{reg})$
at $t\in T(F)_{reg}$ equals the orbital integral
$F_{\overline{f}}\in C_c^\infty(\overline{T(F)}_{reg})_{gen}$
at $(t,1)\in \overline{T(F)}_{reg}$.
\end{definition}

\subsection{Corresponding spherical functions, $p> 2$}

We sketch the Satake correspondence for $\overline{G(F)}$
following \cite{J}, \S 1.3 and using its notation.

Recall that the genuine characters of
$\overline{A(F)}$ are in 1-1 correspondence with
the characters of $A(F)$. Moreover, from results recalled in
the next chapter we deduce that all
the unramified genuine representations of
$\overline{G(F)}$ belong to the unramified principal series.
Replacing $GL(r)$ by
$SL(2)$ in \cite{J}, \S 1.3, we first observe that the
set $\Pi_{unr.p.s.}(\overline{G(F)})_{gen}$ is
parametrized by the $W$-orbit of unramified
genuine characters of $\overline{A(F)}$. Therefore, we have the following

\begin{lemma} The set $\Pi_{unr.p.s.}(\overline{G(F)})_{gen}$ is
in natural 1-1 correspondence with
$\Pi_{unr.p.s.}(G(F))$.
\end{lemma}

In other words, the ``L-group'' of $\overline{G(F)}$ is the
same as that of $G(F)$ - namely, $PGL(2,{\Bbb{C}})$.

Following now \cite{J}, \S 1.4, we obtain a transform
\begin{equation}
\begin{array}{c}
{\cal{H}}(\overline{G(F)},K)_{gen}\rightarrow
{\cal{A}}_{cl.fcn}(PGL(2,{\Bbb{C}}))\\
f\longmapsto S_0 f(\chi)={\rm tr}I_A(\chi,f).
\end{array}
\end{equation}

This is essentially the inverse Fourier transform of the Satake
transform, which we define next. The Satake transform of a
genuine spherical function is defined as in the non-metaplectic Case:
\begin{equation}
Sf(a)=\delta(a)^{1/2}\int_{N(F)}f(an)dn.
\end{equation}
The relationship between this expression and an orbital integral
is well-known - see for example \cite{C}.

This defines an algebra isomorphism
\begin{equation}
\begin{array}{c}
{\cal{H}}(\overline{G(F)},K)_{gen}\rightarrow
{\cal{H}}(\overline{A(F)},K\cap A(F))^W_{gen}\\
f\longmapsto Sf(a).
\end{array}
\end{equation}

Since
${\cal{H}}(\overline{A(F)},K\cap A(F))^W_{gen}\cong
{\cal{H}}(A(F),K\cap A(F))^W$, we have that
\[
{\cal{H}}(\overline{G(F)},K)_{gen}\cong {\cal{H}}(G(F),K).
\]

This implies the matching of orbital integrals for spherical functions.

We shall now make a remark on the matching of
weighted orbital integrals of spherical functions.
Every global smooth genuine function
$\overline{f}=\otimes_v\overline{f}_v$ is a restricted direct product
of local smooth functions, almost all of which are units of the
local Hecke algebra of genuine spherical functions on
$\overline{G(k_v)}$. In order to match weighted orbital integrals
of such global smooth functions, we must match them locally.
A direct calculation of the weighted orbital integrals
shows that if $\overline{f}$ is 
the unit element of the Hecke algebra of genuine spherical functions on
$\overline{G(F)}$ and if $f$ is the
unit element of the Hecke algebra of spherical functions on
$G(F)$ then the weighted orbital integral of $\overline{f}$ on
$\overline{A(F)}$ matches the weighted orbital integral of $f$ on
$A(F)$. We shall use this later in our comparisons in chapter 10.

\section{Weighted characters}

Certain weighted characters occur naturally in the invariant trace formula.
We will need some of their properties to put the trace formula in stable
form. 

\subsection{Irreducible representations of $\overline{A(F)}$}

We use Clifford-Mackey theory to determine the genuine irreducible
representations\ of the local covers $\overline{A(F)}$. We assume that all
representations are smooth (the definition of smooth is recalled in the next
section).

Let $\chi $ denote a genuine (smooth, one-dimensional) character of
$\overline{C}$, where $\overline{C}$ is a maximal abelian subgroup
of $\overline{A(F)}$ as in the previous subsection. Define
$\chi^{a_2}(a_1)=\chi (a_2a_1a_2^{-1})$, for
$a_1\in \overline{C},\;a_2\in \overline{A(F)}$.
The conjugates $\chi ^{a_2}$ are of the form 
\begin{eqnarray}
\chi ^{a_2}(a_1) &=&(t_1,t_2)_n^2\cdot \chi (a_1)\;\;\;\;\;\;\;\;(=\chi
(a_1),\;{\rm if\;}\;a_i\in \overline{C},\;i=1,2)  \nonumber \\
a_1 &=&(\left( 
\begin{array}{cc}
t_1 & 0 \\ 
0 & t_1^{-1}
\end{array}
\right) ,\varsigma _1)\in \overline{C},\;a_2=(\left( 
\begin{array}{cc}
t_2 & 0 \\ 
0 & t_2^{-1}
\end{array}
\right) ,\varsigma _2)\in \overline{A(F)}.
\end{eqnarray}
In particular, there are no genuine 1-dimensional representations of $%
\overline{A(F)}$ unless $n=1,2$.

\begin{lemma}
The characters of $A(F)/C$ are in 1-1 correspondence
with $A(F)/C$ itself. Explicitly, the correspondence is given by
$x\longmapsto \chi _x$, where $\chi _x(y)=(x,y)_n^2$.
\end{lemma}
proof: Pontryagin duality. $\Box$

Therefore, there are $\#[A(F)/C$ $]$ distinct characters of
$A(F)$ which are trivial on $C$. The
number of distinct characters of $C$ of the form $\chi _x\;(x\in A(F))$ is
therefore $\#[A(F)/C$ $]$. This proves the following

\begin{lemma}
Let $\chi $ denote a genuine character of $\overline{C}$, where $C$ is as in
the previous subsection. The number of distinct conjugates $\chi ^h$ of
$\chi $\ is $\#[A(F)/C]$.
\end{lemma}

\begin{lemma}
(Clifford-Mackey) If $\chi \in \Pi (\overline{C})$ then
$\chi ^h\neq \chi $ for all $h\in \overline{A(F)}/\overline{C}$ and 
\[
Ind_{\overline{C}}^{\overline{A(F)}}(\chi )\cong
Ind_{\overline{C}}^{\overline{A(F)}}(\chi ^h),
\;\;\;h\in \overline{A(F)},
\]
is both irreducible and genuine.
\end{lemma}

\begin{remark}
This would follow from an irreducibility criterion of Mackey (stated in
\cite{Se}, section 7.4 for finite groups) if we knew it applied in
our case. To see
that the proof in \cite{Se} extends to the case
considered here we must verify
the Frobenius reciprocity theorem in our case. (As noted in
\cite{C}, Frobenius
reciprocity is false for totally disconnected groups in general.) 
\end{remark}

{\sf proof:} In this proof, let $H<G$ denote
two groups of td type in the sense of \cite{C}, section 1,
such that $H$ is finite index in $G$. Then the inclusion map 
$\varphi :H\hookrightarrow G$ is open. Let $W$ be a smooth representation of 
$H$ and let $V$ be a smooth representation of $G$. By \cite{C}, 1.7(c),
Frobenuis reciprocity holds in this case and we have
\begin{equation}
\dim Hom_H(W,V|_H)=\dim Hom_G(Ind_H^GW,V).
\end{equation}

We apply this to the proof of \cite{Se}, Corollary in section 7.4
with $H=\overline{C}$ and $G=\overline{A(F)}$. 

The representation $V=Ind_{\overline{C}}^{\overline{A(F)}}(\chi )$ is
irreducible if and only if
\[\dim Hom_{\overline{A(F)}}(Ind_{\overline{C}}^{\overline{A(F)}}(\chi ),
Ind_{\overline{C}}^{\overline{A(F)}}(\chi ))=1.
\]
Since 
\begin{equation}
Ind_{\overline{C}}^{\overline{A(F)}}(\chi )|_{\overline{C}}=\sum_{h\in 
\overline{A(F)}/\overline{C}}\chi ^h,
\end{equation}
we have $\dim Hom_{\overline{C}}(\chi ,Ind_{\overline{C}}^{\overline{A(F)}%
}(\chi )|_{\overline{C}})=1$ if and only if all the $\chi ^h$ are distinct. $%
\Box $

\begin{lemma}
Let $\sigma \in \Pi (\overline{A(F)})$. Then

(a) 
\[
\sigma |_{\overline{C}}=\bigoplus_{i=1}^mm\chi _i
\]
for some distinct characters $\ \chi _i$ of $\overline{C}$ and some integers 
$m>0$ (all equal).

(b) If $\ \sigma {\rm \;},\sigma ^{\prime }\in \Pi (\overline{A(F)})$
satisfy 
\[
\sigma |_{\overline{C}}=\bigoplus_{i=1}^m\chi _i,
\;\;\;\;\sigma ^{\prime }|_{\overline{C}}
=\bigoplus_{i=1}^{m^{\prime }}\chi _i^{\prime },\;\;
\]
and $\ \chi _i=\ \chi _j^{\prime }$ for some $i,j$ then $m=m^{\prime }$,
$\{\ \chi _i\;|\;1\leq i\leq m\}=\{\ \chi _i^{\prime }\;|\;1\leq i\leq
m^{\prime }\}$, and
$\ \sigma ^{\prime }\cong \mu \otimes \ \sigma {\rm \;}$
for some character $\mu $ of $\overline{A(F)}$ which is trivial
on $\overline{C}$.

(c) $\ \sigma {\rm \;}$ is a subrepresentation of $Ind_{\overline{C}}^{%
\overline{A(F)}}\ \chi $ for some character $\ \chi $ of $\overline{C}$.
\end{lemma}

proof: These all follow from the lemmas in \cite{GeKn}, section 2, except
possibly (c), which follows from Frobenius reciprocity. $\Box $

\begin{corollary}
If $\ \sigma {\rm \;}=Ind_{\overline{C}}^{\overline{A(F)}}\ \chi
\in \Pi (\overline{A(F)})$ then 
\[
{\rm tr}\;\sigma {\rm \;}(a)=\left\{ 
\begin{array}{c}
\sum_{h\in \overline{A(F)}/\overline{C}}\chi
^h(a),\;\;\;\;\;\;\;\;\;\;\;\;\;a\in \overline{C}, \\ 
0,\;\;\;\;\;\;\;\;\;a\in \overline{A(F)}-\overline{C}.
\end{array}
\right. 
\]
Indeed, by equation (5.15) we have that
\[
\sum_{h\in \overline{A(F)}/\overline{C}}\chi^h(a),
\]
vanishes unless $a\in \overline{A(F)^N}$. Therefore,
${\rm tr}\;\sigma {\rm \;}(a)$ is supported on
$\overline{A(F)^N}$.
\end{corollary}

\subsection{Principal series representations}
From results recalled in the next chapter, we have

\begin{lemma}
Assume $\chi =Ind_{\overline{C}}^{\overline{A(F)}}\mu
\in \Pi (\overline{A(F)})$ and suppose that $\chi \neq \chi ^w$.
The character of the
irreducible representation $\pi =I_{A}(\chi )$\ is given on the regular
set by
\[
\Theta _\pi (x)=\left\{
\begin{array}{c}
\frac{tr\chi (a)+tr(\chi (a^w))}{|a_1-a_2|},
\;\;x=a^g,\;{\rm for\;some}\;g\in \overline{G(F)},
\;a\in \overline{C}-\{\pm 1\}, \\
0,\;\;\;\;\;\;\;\;\;\;\;\;\;\;\;\;\;\;\;\;\;\;\;\;\;\;\;\;\;
{\rm \;if\;}\;\;\;
x\in \overline{G(F)}^{\prime }-\overline{C}^{\overline{G(F)}},
\end{array}
\right.
\]
where
\[
\overline{C}^{\overline{G(F)}}=\{a^g\;|\;a\in \overline{C},\;g\in
\overline{G(F)}\}.\
\]
\end{lemma}

\begin{remark} As noted above, on the regular set,
this character vanishes off the conjugacy class of $\overline{A(F)^N}$.
\end{remark}

Now we describe the character of a principal series representation
as the Fourier transform of the function $a\longmapsto J_G(a,f)$.
For the existence of characters of metaplectic covering groups,
see the next chapter.

Let $\chi =Ind_{\overline{C}}^{\overline{A(F)}}\mu \in
\Pi (\overline{A(F)})$,
$\pi =I_{A}(\chi )\in \Pi (\overline{G(F)})$. Then
\begin{equation}
{\rm tr}\chi (a,1)=
\#( \overline{A(F)}/\overline{C})\mu (a,1),
\end{equation}
if $a\in A(F)^N$. This lemma and the Weyl integration formula
imply that, for $f\in C_c^\infty
(\overline{G(F)})$,
\begin{equation}
\begin{array}{c}
\Theta _\pi (f)=\int_{\overline{G(F)}}\Theta _\pi (x)f(x)dx \\
=\frac 12\int_{A(F)}({\rm tr}\chi (a,1)+{\rm tr}(\chi
(a^w,1)))F_f^{A(F)}(a,1)da \\
=\#( \overline{A(F)}/\overline{C})
\int_{A(F)^N}\mu (a,1)F_f^{A(F)}(a,1)da \\
=\#(A(F)/C)
(F_f^{A(F)})^{\wedge }(\mu ) ,
\end{array}
\end{equation}
where the Fourier transform is defined by
\begin{equation}
\begin{array}{c}
\varphi ^{\wedge }(\mu )=\int_{A(F)^N}\mu (x)\varphi
(x)dx \\
\ \ \ \ \ \ \mu \in \Pi (\overline{A(F)^N}).
\end{array}
\end{equation}
Note that this formula for the character extends to the
Schwartz space provided $\mu $ (and hence $\pi $) is unitary.

\begin{corollary}
Let $n=2$. If $\phi\in C^\infty_c(\overline{G(F)})_{gen}$
is related to $f\in C_c^\infty(G(F))$ by
\[
\phi=\overline{f},
\]
as in section 5.4 above, and if
$\overline{\pi}\in \Pi(\overline{G(F)})$ is related to
$\pi\in\Pi(G(F))$ by
\[
\overline{\pi}
=I_A(\overline{\chi}),\ \ \
\pi=I_A(\chi),
\]
for $\chi\in\Pi(A(F))$, then
\[
\Theta_{\overline{\pi}}(\phi)=\Theta_{\pi}(f).
\]
\end{corollary}

\begin{lemma} Let $n=2$. For $\pi=I_A(\chi)$,
$\chi=Ind_{\overline{C}}^{\overline{A(F)}}(\mu)$, as above,
the distribution $f\longmapsto \Theta_\pi(f)$ is
invariant under conjugation by elements of
$\overline{G_1(F)}$.
Moreover, its inverse Fourier transform
distribution $f\longmapsto F_f^{\overline{A(F)}}$ is
invariant under conjugation by elements of
$\overline{G_1(F)}$.
\end{lemma}

\begin{lemma}
If $F(\sigma)={\rm tr}(I_A(\sigma,f))$ then
\[
F^\vee(\gamma)
=|a_1-a_2|J_G(\gamma,f),
\]
for $\gamma =(a,\zeta)\in \overline{A(F)}$, $a\not= \pm 1$.
\end{lemma}

proof: This follows from our results above. $\Box$

The following few result will not be needed.

\begin{lemma} If $\theta_1$ is an irreducible genuine representation
of $\overline{T_1({\Bbb Q}_p)}$ and $\pi_1(\theta_1)$ is the
irreducible genuine admissible representation of
$\overline{G_1({\Bbb Q}_p)}$ associated to it in \cite{J}
(at least if $n$ is odd), let
$\theta =\theta_1|_{\overline{T({\Bbb{Q}}_p)}}$ and let
$\pi(\theta)$ denote one of the irreducible
constituents of
$\pi(\theta) =\pi_1(\theta_1)|_{\overline{G({\Bbb{Q}}_p)}}$.
If $n>1$ then
\[
\# L(\pi(\theta))=
\prod_{v|p}\#( k_v^\times/k_v^{\times n}).
\]
\end{lemma}

proof: For the proof in the case $n=2$, see \cite{GPS}, \S 1.
The proof in general is similar, noting that
the Hilbert symbol provides a perfect pairing
$F^\times/F^{\times n}\times F^\times/F^{\times n}
\rightarrow \mu_n$. $\Box$

\begin{lemma}
For each elliptic torus $T$ and each
$\overline{f}\in {\cal{H}}(\overline{G({\Bbb{A}}_S)})_{gen}$,
we have
\[
\begin{array}{c}
<\theta ,I_{*}^\kappa (\overline{f} )> \\
=[D(T/{\Bbb{Q}}):D(T/{\Bbb{A}}_S)]^{-1}
\prod_{p\in S}
\sum_{\delta\in D(T/{\Bbb Q}_p)}
\sum_{\pi _p\in \Pi (\theta_p^\delta)}
\kappa _{*}(\pi _p)\chi _{\pi_p}(f_p),
\end{array}
\]
where $\kappa _{*}(\pi _p):=\kappa (\delta)$ if
$\pi _p\in \Pi(\theta_p^\delta)$,
${\theta }_p^h$ denotes the conjugate of ${\theta }_p$ by $h$,
${\theta }_p$ is a unitary character
of $T({\Bbb{Q}}_p)$, and $f_p(g)=\overline{f}_p(g,1)$.
\end{lemma}

\begin{remark}
This result is essentially \cite{LL}, (2.4-5) and p. 762.
This is, after a suitable
choice of global measures, related to a special case of
\cite{J}, Theorem 3.41 which was only proven under in the case $n$ odd.
\end{remark}

proof: We prove the local analog:
\[
\begin{array}{c}
(I_{*}^\kappa (f ))^\wedge(\theta) \\
=\sum_{\delta\in D(T/{\Bbb Q}_p)}
\sum_{\pi\in L({\pi }({\theta }^\delta))}
\kappa _{*}(\pi)\chi _{\pi }(f).
\end{array}
\]
This is a consequence of the above-mentioned result of
\cite{LL} and Definition 73 above. $\Box$

\subsection{Weighted characters }

The operator
$M(\mu):I_A(\mu)\rightarrow I_A(\mu^w)$ denotes the usual
intertwining operator as in \cite{Ar}, \cite{Mo1}, \cite{Mo2}.
We will eventually
assume that the normalized intertwining operator
$R(\mu)$ satisfies the assumptions in section 4 of chapter 9.
Here, however, we simply assume that there is a scalar valued
function (the normalizing factor) $m(\mu)$ such that
\[
R(\mu)=m(\mu)^{-1}M(\mu).
\]


Let $\delta(g)=|a_1/a_2|$, if $g=(nak,\zeta)\in \overline{G_1(F)}$. Let
$N(g)$ denote the operator on the space of functions on $K$ given by
\[
f(k)\longmapsto f(k)\log\delta(kg).
\]
This operator also occurs in connection with the terms in the
invariant trace formula associated to the normalized intertwining
operators.

\begin{lemma}
For $g\in \overline{G_1(F)}$, we have
\[
{\rm tr}(I_A(\sigma,f)N(g))
=\int_K K_{I_A(\sigma)}(k,k)\log\delta(kg)dk,
\]
here
\[
K_{I_A(\sigma)}(k,k')
\int_{\overline{A(F)}}\int_{N(F)}
\delta(a)^{1/2}{\rm tr}\chi(a)f(k^{-1}ank')dnda.
\]
\end{lemma}

\begin{lemma}
For $g\in \overline{G_1(F)}$, if
$F(\sigma)={\rm tr}(I_A(\sigma,f)N(g))$ then
\[
J_A(a,f^{g^{-1}})-J_A(a,f)
=F^\vee(a)+F^\vee(a^w)
\]
\end{lemma}
proof: We shall follow \cite{LL}, Lemma 3.2, p. 743.
We have
\[
J_A(a,f^{g^{-1}})-J_A(a,f)
=\int_{\overline{G(F)}_a\backslash \overline{G(F)}}
f(x^{-1}ax)(v_M(gx)-v_M(x))dx.
\]
This is, by a lemma above, zero unless
$a\in\overline{A(F)^N}$. We first show that

{\bf Claim} $v_A(nk)=\log\delta(nk)+\log\delta(wnk)$.

This is verified by direct calculations using the matrix calculation
\[
\left(
\begin{array}{cc}
0 & -1\\
1 & 0
\end{array}
\right)
\left(
\begin{array}{cc}
1 & x\\
0 & 1
\end{array}
\right)
=
\left(
\begin{array}{cc}
1 & -x^3\\
0 & 1
\end{array}
\right)
\left(
\begin{array}{cc}
x & 0\\
0 & x^{-1}
\end{array}
\right)
\left(
\begin{array}{cc}
1 & 0\\
x^{-1} & 1
\end{array}
\right).
\]

Since
$\delta(g_1g_2)=\delta(g_1)\delta(g_2)$,
the above claim implies imply that
$v_A(x)=\log\delta(x)+\log\delta(wx)$ for all $x$.
These identities are all we need to carry over the argument
in \cite{LL}, p. 743. $\Box$

The following result will not be needed.

\begin{lemma} For $g\in\overline{G_1(F)}$, we have
\[
\begin{array}{c}
{\rm tr}[R(\mu)^{-1}R'(\mu)I_A(\mu,f^{g^{-1}})]
-{\rm tr}[R(\mu)^{-1}R'(\mu)I_A(\mu,f)]\\
={\rm tr}[I_A(\mu,f)N(g))+
{\rm tr}[I_A(\mu^w,f)N(g)],
\end{array}
\]
where $N(g)$ is defined above.
\end{lemma}

For the proof in the case n=1, see \cite{LL}, Lemma 3.6.
The proof for $n>1$ is the same, although we include a sketch of
it for the sake of completeness.
proof: First, we show that the Lemma follows from the identity
\begin{equation}
\pi_{{\mu}^w}(g)R'({\mu})I_A(\mu)(g)^{-1}-R'(\mu)
=R(\mu)N(g)+N(g)R(\mu).
\end{equation}
Indeed, this equation yields
\[
\begin{array}{c}
I_A(\mu^w)(g)R'({\mu})I_A(\mu)(g)^{-1}
I_A(\mu,f)-R'(\mu)I_A(\mu,f)\\
=R(\mu)N(g)I_A(\mu,f)+
N(g)R(\mu)I_A(\mu,f),
\end{array}
\]
and this gives
\[
\begin{array}{c}
R({\mu})^{-1}I_A(\mu^w)(g)R'({\mu})
I_A(\mu)(g)^{-1}I_A(\mu,f)
-R(\mu)^{-1}R'(\mu)I_A(\mu,f)\\
=I_A(\mu)(g)[R(\mu)^{-1}R'({\mu})
I_A(\mu)(g)^{-1}
I_A(\mu,f)I_A(\mu)(g)]I_A(\mu)(g)^{-1}
-R(\mu)^{-1}R'(\mu)I_A(\mu,f)\\
=N(g)I_A(\mu,f)
+R(\mu)^{-1}N(g)R(\mu)I_A(\mu,f)\\
=N(g)I_A(\mu,f)+R(\mu)^{-1}
[N(g)I_A(\mu^w,f)]R(\mu).
\end{array}
\]
This implies the Lemma.
\par
Now we verify (5.21). By the metaplectic Iwasawa decomposition, we
may write $wnk=n_1a_1k_1$, for some $n_1\in N(F)$,
$a_1\in\overline{A(F)}$,$k_1\in\overline{K}$ and a unique $p(a_1)\in A(F)$.
For $k\in\overline{K}$ and $\phi\in V(\pi_s)$, we have
$$
M'(\mu)\phi (k):=\int_{N(F)}{d\over ds}\phi (wnk)dn=
\int_{N(F)}\log \delta(p(a_1))\phi (wnk)dn.
$$
Let
\[
\begin{array}{c}
kg=n_2a_2k_2,\\
wnk_2=n_1'a_1'k_1',\\
k_1g=n_3a_3k_3,
\end{array}
\]
so that
\[
\begin{array}{c}
\pi_{{\mu}^w}(g)M'(\mu)\pi_{{\mu}}(g)^{-1}\phi (k)
=\int_{N(F)}\log \delta(p(a_1'))\phi (wnk)dn,\\
\log \delta (p(a_1'))
=\log \delta (p(a_1))+\log \delta (p(a_2))+\log \delta (p(a_3)),\\
M(\mu)N(g)\phi (k)
=\int_{N(F)}\log \delta(p(a_2))\phi (wnk)dn,\\
N(g)M(\mu)\phi (k)
=\int_{N(F)}\log \delta(p(a_3))\phi (wnk)dn.
\end{array}
\]
From these identities, the Lemma follows as in \cite{LL}, p. 746.
$\Box$

\subsection{A matching of weighted characters}

\begin{lemma}
Let $n=2$. If $\overline{f}\in C^\infty_c(\overline{G(F)})_{gen}$
is related to $f\in C_c^\infty(G(F))$ as in section 5.2
above. Then
\[
{\rm tr}(R(\overline{\chi})^{-1}R'(\overline{\chi})
I_A(\overline{\chi},\overline{f}))=
{\rm tr}(R(\chi)^{-1}R'(\chi)I_A(\chi,f)).
\]
\end{lemma}

\begin{remark} In spite of its apparent simplicity, the proof
below implicitly uses the assumptions in chapter 9, \S 4
and results of \cite{A6}.
\end{remark}

proof: This is equivalent to saying that the Fourier transforms of both
sides (with respect to $\chi$) are the same at $a\in A(F)_{reg}$,
resp. $(a,1)\in \overline{A(F)}_{reg}$.
By (9.68) below (which is a special case of
\cite{A6}, Proposition 4.1), this is in turn equivalent to saying that
\begin{equation}
(I_{\cal{O}}^A)^\wedge(\phi_A^G(f))
=(I_{\overline{\cal{O}}}^{\overline{A}})^\wedge
(\phi_{\overline{A}}^{\overline{G}}(\overline{f}))
\end{equation}
where we have distinguished objects on $G(F)$ and their analogs on
$\overline{G(F)}$ by an overline. We have already shown that
the distribution on $G(F)$
\[
f\longmapsto J(f)=J_{\cal{O}}(f)-
J_{\overline{\cal{O}}}(\overline{f}))
\]
is trivial. By \cite{A6}, Proposition 4.1, there are {\em unique}
invariant distributions
$I_{\cal{O}}^{M}$, $I_{\overline{\cal{O}}}^{\overline{M}}$ such that
\[
J(f)=\sum_M [ c(M)\hat{I}_{\cal{O}}^M(\phi_M^G(f))-
 c(\overline{M})\hat{I}_{\overline{\cal{O}}}^{\overline{M}}
(\phi_{\overline{M}}^{\overline{G}}(\overline{f})).
\]
By uniqueness,
\[
\hat{I}_{\cal{O}}^M(\phi_M^G(f))
={c(\overline{M})\over c(M)}
\hat{I}_{\overline{\cal{O}}}^M
(\phi_{\overline{M}}^{\overline{G}}(\overline{f})).
\]
We see that $c(A)=c(\overline{A})=1/2$ by (9.68) below
and its non-metaplectic analog, so the result follows. $\Box$



